% source: erdos-872/current_state.md (Round 7 / Round 8 upper bounds)
% source: erdos-872/researcher-08-5-16-improvement.md
% source: erdos-872/lean/README.md

\section{Intermediate Upper Bounds}\label{sec:intermediate-upper}

Before turning to the Round~15 piecewise-density argument, we record the two
earlier upper bounds obtained from the odd-prime-prefix Shortener strategy.  The
same basic mechanism appears in both arguments:
\begin{enumerate}
  \item Shortener plays a controlled prefix of legal odd primes;
  \item the odd-part map injects the final antichain into the odd integers
    avoiding those primes;
  \item a low-order Bonferroni truncation bounds the surviving odd integers.
\end{enumerate}

\subsection{The \texorpdfstring{$13/36$}{13/36} bound}

\begin{theorem}\label{thm:1336}
Under the strategy ``play the smallest legal odd prime,'' Shortener has
\[
  \L(n) \le \left(\frac{13}{36}+o(1)\right)n.
\]
\end{theorem}

\begin{proof}[Proof sketch]
Let $q_1<q_2<\cdots$ be the odd primes played by Shortener.  The key uniform
input is the Chebyshev-type bound
\[
  q_j \le \left(\frac32+o(1)\right)j\log n,
\]
which guarantees enough legal odd primes during the prefix.

Choose $t$ minimal with
\[
  s_t := \sum_{j\le t}\frac1{q_j} \ge \frac13-o(1).
\]
Because every $q_j \ge 3$, minimality forces
\[
  \frac13-o(1) \le s_t \le \frac23.
\]
Let $D=\{q_1,\dots,q_t\}$.  The odd-part map sends the final antichain
injectively into the set of odd integers not divisible by any prime in $D$,
together with the already-played prefix itself.  If $N_D(n)$ denotes the number
of odd integers up to $n$ coprime to $D$, then
\[
  \L(n) \le t + N_D(n).
\]
Since $t=o(n)$, it remains to bound $N_D(n)$.

Second-order Bonferroni on the odd integers gives
\[
  N_D(n)
  \le
  \frac n2
  \left(1-s_t+\frac{s_t^2}{2}\right)
  + o(n).
\]
The function $f(s)=1-s+s^2/2$ is decreasing on $[0,1]$, and the truncation
ensures $s_t \in [1/3-o(1),2/3] \subset [0,1]$.  Hence
\[
  f(s_t) \le f(1/3-o(1)) = \frac{13}{18}+o(1).
\]
Therefore
\[
  \L(n) \le \frac n2 \cdot \frac{13}{18} + o(n)
  = \left(\frac{13}{36}+o(1)\right)n.
\]
\end{proof}

\begin{remark}
This theorem is fully verified in Lean with no remaining sorries.  The
formalized artifact uses a slightly simplified input set in order to avoid a
Chebyshev bottleneck, but the coefficient is the same.
\end{remark}

\subsection{The \texorpdfstring{$5/16$}{5/16} bound}

The next observation is that the same odd-prime-prefix idea admits a much longer
prefix than the original $13/36$ proof used.

\begin{theorem}\label{thm:516}
For every fixed $A>2$, the odd-prime-prefix Shortener strategy gives
\[
  \L(n) \le
  \left(
    \frac12 - \frac{1}{2A} + \frac{1}{4A^2} + o(1)
  \right)n.
\]
Letting $A \downarrow 2$ yields
\[
  \L(n) \le \left(\frac{5}{16}+o(1)\right)n.
\]
\end{theorem}

\begin{proof}[Proof sketch]
Set
\[
  k = \left\lfloor \frac{n}{2A\log n}\right\rfloor,
\]
and let Shortener play the smallest legal odd prime on her first $k$ turns.
If the game ends earlier, then the game length is already $o(n)$, so assume the
prefix exists and denote it by $q_1<\cdots<q_k$.

The longer prefix is available because the same argument as above now gives
\[
  q_j \le A j \log n \qquad (1\le j\le k).
\]
Consequently
\[
  \sum_{j\le k}\frac1{q_j}
  \ge
  \frac{1}{A\log n}\sum_{j\le k}\frac1j
  =
  \frac1A+o(1).
\]
Choose $t\le k$ minimal with
\[
  s_t := \sum_{j\le t}\frac1{q_j} \ge \frac1A-o(1),
\]
and again let $D=\{q_1,\dots,q_t\}$.  The same odd-part injection and
second-order Bonferroni argument give
\[
  \L(n) \le t + \frac n2\left(1-s_t+\frac{s_t^2}{2}\right) + o(n).
\]
Here $t=o(n)$, and $f(s)=1-s+s^2/2$ is decreasing on $[0,1]$, so
\[
  \L(n) \le
  \frac n2\left(1-\frac1A+\frac{1}{2A^2}\right)+o(n)
  =
  \left(\frac12 - \frac{1}{2A} + \frac{1}{4A^2} + o(1)\right)n.
\]
Sending $A\downarrow 2$ gives the bound $5/16$.
\end{proof}

\begin{remark}
The one remaining Lean gap in the current artifact is Chebyshev-adjacent and
does not affect the mathematical status of the theorem.  For the paper, the
important point is conceptual: the `13/36' proof did not exhaust the
odd-prime-prefix method.  The correct limiting constant for that family is
$5/16$.
\end{remark}

\subsection{Why this family stops at \texorpdfstring{$5/16$}{5/16}}

The odd-prime-prefix strategy is elegant but not universal.  Later obstruction
work shows that Prolonger can force
\[
  \sum_{j} \frac1{q_j} \le \frac12 + o(1)
\]
against this entire style of argument.  Since the second-order truncation
coefficient is
\[
  \frac12\left(1-s+\frac{s^2}{2}\right),
\]
the best constant accessible to this family is exactly $5/16$, achieved in the
limit $s=1/2$.

This is why the main theorem below $0.19n$ must leave the odd-prime-prefix
framework and instead exploit the Round~15 piecewise density of the actual prime
sequence.
